%!TEX program = xelatex
\documentclass[11pt,a4paper]{article}

\usepackage{fontspec}
\usepackage{polyglossia}
\setmainlanguage{vietnamese}
\usepackage[margin=1in]{geometry}
\usepackage{amsmath,amssymb}
\usepackage{float}
\usepackage{booktabs}
\usepackage{array}
\usepackage{graphicx}
\usepackage{enumitem}
\usepackage{xcolor}
\usepackage[hidelinks]{hyperref}

\title{Hệ thống RAG Chẩn đoán Bệnh dựa trên Triệu chứng\\
       \large Phần 8 --- Ablation Study và Phân tích Bổ sung}
\author{AIO Vietnam 2026 --- Module 1 Project\\
        \small Phụ trách Phần 8 \& QA: Nguyễn Thị Diễm Hương}
\date{\today}

\begin{document}
\maketitle

\section{Kết quả thực nghiệm (bối cảnh cho Phần 8)}
\label{sec:ketqua}

Phần này tóm tắt các kết quả chính được Phần~\ref{sec:ablation} sử dụng làm
điểm tham chiếu. Toàn bộ đánh giá thực hiện trên bộ dữ liệu
DDXPlus~\cite{ddxplus2022} với $n=132{,}448$ truy vấn triệu chứng đã chuẩn hoá.

\begin{table}[H]
\centering
\small
\caption{Kết quả truy xuất offline (Python, $n=132{,}448$). Hai biến thể trường
văn bản: \texttt{keyword\_text} và \texttt{keyword\_text+description}.}
\label{tab:offline}
\begin{tabular}{lcccc}
\hline
\textbf{Phương pháp} & \multicolumn{2}{c}{\textbf{Hit@1 (\%)}} &
\multicolumn{2}{c}{\textbf{MRR}}\\
 & \texttt{kw\_text} & \texttt{+desc} & \texttt{kw\_text} & \texttt{+desc}\\
\hline
Dense (bge-small-en-v1.5) & 83.11 & 84.28 & 0.8790 & 0.8902\\
Hybrid-RRF ($k=60$)       & 91.17 & 91.95 & 0.9385 & 0.9441\\
BM25                      & 98.47 & 98.78 & 0.9921 & 0.9937\\
\hline
\end{tabular}
\end{table}

\begin{figure}[H]
\centering
\fbox{\parbox[c][3.2cm][c]{0.7\linewidth}{\centering\small
\textit{Biểu đồ cột Hit@1 theo phương pháp}\\[4pt]
BM25 $98.78\% >$ Hybrid-RRF $91.95\% >$ Dense $84.28\%$\\[2pt]
(biến thể \texttt{keyword\_text+description})}}
\caption{Hit@1 theo phương pháp truy xuất (offline, biến thể có
\texttt{description}). Thứ hạng BM25 $>$ Hybrid-RRF $>$ Dense ổn định ở mọi cấu
hình.}
\label{fig:bar}
\end{figure}

\begin{table}[H]
\centering
\small
\caption{Đối chiếu Hit@1 (\%) giữa hạ tầng offline (Python) và live (OpenSearch),
cùng KB / normalizer / ground truth. Thứ hạng phương pháp được bảo toàn.}
\label{tab:livevsoffline}
\begin{tabular}{lccc}
\hline
\textbf{Phương pháp} & \textbf{Offline} & \textbf{Live (OpenSearch)} &
\textbf{$\Delta$}\\
\hline
BM25        & 98.78 & 98.65 & $-0.13$\\
Hybrid-RRF  & 91.95 & 92.30 & $+0.35$\\
Dense       & 84.28 & 79.54 & $-4.74$\,$^{\dagger}$\\
\hline
\end{tabular}\\[2pt]
{\footnotesize $^{\dagger}$ Nhánh Dense live dùng biến thể query-embed legacy nên
thấp hơn offline; ghi rõ để tránh hiểu nhầm là mâu thuẫn số liệu.}
\end{table}

\begin{table}[H]
\centering
\small
\caption{Bản chạy đầy đủ trên OpenSearch (live, $n=132{,}448$, ngày 26/06,
seed mặc định). Δ so với mốc offline biến thể \texttt{+description}.}
\label{tab:livefull}
\begin{tabular}{lccc}
\hline
\textbf{Phương pháp} & \textbf{Hit@1 (\%)} & \textbf{Offline (\%)} &
\textbf{$\Delta$ (điểm \%)}\\
\hline
BM25        & 98.65 & 98.78 & $-0.13$\\
Hybrid-RRF  & 92.30 & 91.95 & $+0.35$\\
Dense       & 79.54 & 84.28 & $-4.74$\\
\hline
\end{tabular}
\end{table}

\clearpage

\section{Ablation Study và phân tích bổ sung}
\label{sec:ablation}

Mục này tách riêng đóng góp của từng thành phần trong pipeline truy xuất nhằm
trả lời câu hỏi: \emph{mỗi lựa chọn thiết kế đem lại bao nhiêu giá trị?} Trừ khi
ghi chú khác, Hit@1 trích từ Bảng~\ref{tab:offline} (offline, $n=132{,}448$).

\subsection*{(A) Ảnh hưởng của trường \texttt{description} (A/B test)}
Bật/tắt \texttt{description} trong cả \texttt{keyword\_text} lẫn embedding cho cải
thiện Hit@1 \emph{nhất quán} ở cả ba phương pháp:
Dense $83.11 \rightarrow 84.28\%$ (\textbf{+1.17} điểm \%),
Hybrid-RRF $91.17 \rightarrow 91.95\%$ (\textbf{+0.78}),
BM25 $98.47 \rightarrow 98.78\%$ (\textbf{+0.31}).
Mức lợi giảm dần theo độ "đã bão hoà" của phương pháp: Dense hưởng lợi nhiều nhất
vì embedding khai thác được ngữ cảnh ngữ nghĩa bổ sung; BM25 vốn gần trần nên lợi
ít. \textbf{Kết luận:} giữ \texttt{description} ở cả hai biểu diễn.

\subsection*{(B) Ảnh hưởng của chiến lược truy xuất}
Thứ tự Hit@1 ổn định \textbf{BM25 $>$ Hybrid-RRF $>$ Dense} ở cả hai hạ tầng.
Đáng chú ý, Hybrid-RRF cao hơn Dense thuần \emph{chỉ nhờ} thành phần BM25 bên
trong RRF, nhưng vẫn \textbf{thấp hơn BM25 đơn lẻ}: ghép thêm nhánh Dense vào lại
\emph{kéo Hit@1 xuống}. Nguyên nhân là đặc thù DDXPlus -- truy vấn là token triệu
chứng đã chuẩn hoá, khớp từ khoá gần như hoàn hảo, nên tín hiệu lexical rất mạnh
và embedding ngữ nghĩa làm "nhoè" tín hiệu đó. Đây là phát hiện ngược trực giác
cần nhấn mạnh: trên dạng dữ liệu này, hybrid \emph{chưa} mang lại lợi ích như kỳ
vọng lý thuyết.

\subsection*{(C) Tính nhất quán hạ tầng (offline Python vs OpenSearch live)}
Đối chiếu cùng KB, cùng normalizer, cùng ground truth, chỉ khác hạ tầng: chênh
lệch Hit@1 nhỏ và \textbf{thứ hạng phương pháp được bảo toàn}
(Bảng~\ref{tab:livevsoffline}). Bản chạy đầy đủ trên OpenSearch ($n=132{,}448$)
cho BM25 $98.65\%$, Hybrid-RRF $92.30\%$, Dense $79.54\%$ -- thống nhất với mốc
offline ở mức điểm phần trăm và giữ nguyên thứ hạng. (Nhánh Dense live thấp hơn
offline do dùng biến thể query-embed legacy; ghi rõ để tránh hiểu nhầm là mâu
thuẫn số liệu.)

\subsection*{(D) An toàn lâm sàng: tỉ lệ bỏ sót bệnh nặng (MSR)}
Ngoài Hit@1 trung bình, một hệ chẩn đoán cần được đánh giá theo \emph{rủi ro bỏ
sót bệnh nặng}. Định nghĩa \textbf{MSR (Missed-Severe Ratio)} là tỉ lệ ca thuộc
nhóm bệnh có severity $\le 2$ (1--2 là nặng nhất) mà nhãn gold \emph{không} nằm ở
top-1. Tính trực tiếp trên BM25 ($36{,}774$ ca nặng):
\texttt{keyword\_text} cho MSR $=1.57\%$ ($576/36{,}774$),
\texttt{keyword\_text+description} cho MSR $=1.19\%$ ($439/36{,}774$).
Trường \texttt{description} không chỉ tăng Hit@1 mà còn \textbf{giảm bỏ sót bệnh
nặng $\approx 24\%$ tương đối}. Toàn bộ ca bỏ sót tập trung ở nhóm tim mạch dễ
nhầm lẫn (Unstable angina, Myocarditis, Stable angina), gợi ý hướng ưu tiên cải
thiện ở Phase~2 (Bảng~\ref{tab:msr}).

\begin{table}[H]
\centering
\small
\caption{MSR (severity $\le 2$) trên BM25, $n_{\text{severe}}=36{,}774$. Nguồn:
\texttt{exp02\_per\_disease.csv}.}
\label{tab:msr}
\begin{tabular}{lccc}
\hline
\textbf{Biến thể BM25} & \textbf{Ca nặng bỏ sót} & \textbf{MSR} & \textbf{Hit@1}\\
\hline
\texttt{keyword\_text}              & 576 & 1.57\% & 98.47\%\\
\texttt{keyword\_text+description}  & 439 & 1.19\% & 98.78\%\\
\hline
\end{tabular}
\end{table}

\subsection*{(E) Độ tin cậy thống kê (Bootstrap 95\% CI)}
Để khẳng định chênh lệch giữa hai biến thể không phải nhiễu, áp dụng bootstrap
($B=1000$, seed $42$) cho Hit@1 trên toàn $n=132{,}448$:
\texttt{keyword\_text} $98.47\%$, CI $95\%=[98.40,\,98.54]\%$;
\texttt{keyword\_text+description} $98.78\%$, CI $95\%=[98.72,\,98.84]\%$.
Hai khoảng tin cậy \textbf{không giao nhau}, xác nhận lợi ích của
\texttt{description} là có ý nghĩa thống kê chứ không do dao động lấy mẫu.

\subsection*{(F) Phân tích dự kiến và tham số RRF (Phase 2)}
Các hướng ablate tiếp theo: (i) \textbf{re-ranking} hai tầng với
\texttt{bge-reranker-base} -- biên độ cải thiện ở Phase~1 kỳ vọng rất nhỏ vì BM25
đã gần trần, giá trị thực rõ hơn trên truy vấn ngôn ngữ tự nhiên;
(ii) \textbf{tinh chỉnh RRF} -- tăng tỉ trọng nhánh BM25 hoặc tăng
\texttt{rank\_constant} ($k=60$ hiện tại); (iii) \textbf{cascade}
BM25 $\rightarrow$ Dense chỉ khi top-1 thiếu tự tin; (iv) thay embedding y khoa
chuyên biệt cho \texttt{bge-small-en-v1.5} tổng quát. Mục tiêu là kiểm chứng giá
trị kiến trúc hai tầng trên truy vấn ngôn ngữ tự nhiên thay vì token chuẩn hoá.

\subsection*{(G) Bảng chỉ số QA bổ sung}
Để khép lại phần phân tích, Bảng~\ref{tab:qa} tổng hợp các chỉ số QA bổ sung. Hai
chỉ số NDCG@5 và ICD-10 chapter/block Hit@1 cần dữ liệu dự đoán per-case (bệnh
top-$k$ ở mỗi ca) -- file summary hiện chỉ tổng hợp per-disease nên tạm để trống
kèm công thức tính sẵn cho Phase~2.

\begin{table}[H]
\centering
\small
\caption{Chỉ số QA bổ sung (Người tính: Nguyễn Thị Diễm Hương). MSR và CI tính
trên BM25 \texttt{keyword\_text+description}, $n=132{,}448$.}
\label{tab:qa}
\begin{tabular}{p{5.6cm}p{3.4cm}p{3.2cm}}
\hline
\textbf{Chỉ số QA} & \textbf{Ý nghĩa} & \textbf{Giá trị}\\
\hline
NDCG@5 (graded relevance) & Chất lượng xếp hạng có trọng số & \textit{[cần per-case rank]}\\
Tier-1.5: ICD-10 chapter Hit@1 & Đúng chương ICD-10 & \textit{[cần dự đoán/ca]}\\
Tier-1.5: ICD-10 block Hit@1 & Đúng block ICD-10 & \textit{[cần dự đoán/ca]}\\
MSR (Missed-Severe Ratio, severity $\le 2$) & Tỉ lệ bỏ sót bệnh nặng & \textbf{1.19\%} (439/36{,}774)\\
Bootstrap 95\% CI (B=1000, seed 42) & Khoảng tin cậy Hit@1 & $[98.72,\,98.84]\%$\\
\hline
\end{tabular}
\end{table}

\begin{thebibliography}{9}
\bibitem{ddxplus2022}
Arsene Fansi Tchango, Rishab Goel, Zhi Wen, Julien Martel, Joumana Ghosn.
\emph{DDXPlus: A New Dataset for Automatic Medical Diagnosis}.
Advances in Neural Information Processing Systems (NeurIPS), Datasets and
Benchmarks Track, 2022.
\end{thebibliography}

\end{document}